\noindent
{\Large \textbf{\section{Постановка задачи}}}\par
Цель работы: реализовать калькулятор «большой» конечной арифметики $<Z_{11}^8;+,*>$ для варианта D2, который задан правилом «+1»: \par 

\vspace{8pt}

\begin{center}
    \begin{tabular}{|c|c|c|c|c|c|c|c|c|c|c|c|c|}
    \hline
    \multicolumn{2}{|c|}{Вариант} & \multicolumn{11}{c|}{Правило «+1»} \\
    \hline
    № & $Z_i$ & a & b & c & d & e & f & g & h & i & j & k \\
    \hline
    D2 & $Z_{11}$ & b & d & a & f & c & i & j & e & k & h & g \\
    \hline
    \end{tabular}
\end{center}

\vspace{8pt}

\noindent и выполняются свойства коммутативности $(+,*)$, ассоциативности $(+, *)$,
дистрибутивности * относительно +, заданы аддитивная единица «a» и
мультипликативная единица «b», а также выполняется свойство: для ($\forall$х) х*a=a.

\vspace{10pt}
Для достижения этой цели необходимо решить задачи:
%\vspace{7pt}
\begin{enumerate}[itemsep=-3pt]
    \item Создать сложение в «малой» конечной арифметике в соответствии с правилом +1.
    \item На основе сложения в «малой» конечной арифметике создать умножение, вычитание, деление.
    \item На основе «малой» конечной арифметики реализовать соответствующие действия для «большой» конечной арифметики.
    \item Создать методы для обработки ввода отрицательных чисел.
\end{enumerate}

\newpage